% Options for packages loaded elsewhere
\PassOptionsToPackage{unicode}{hyperref}
\PassOptionsToPackage{hyphens}{url}
\documentclass[
]{article}
\usepackage{xcolor}
\usepackage{amsmath,amssymb}
\setcounter{secnumdepth}{-\maxdimen} % remove section numbering
\usepackage{iftex}
\ifPDFTeX
  \usepackage[T1]{fontenc}
  \usepackage[utf8]{inputenc}
  \usepackage{textcomp} % provide euro and other symbols
\else % if luatex or xetex
  \usepackage{unicode-math} % this also loads fontspec
  \defaultfontfeatures{Scale=MatchLowercase}
  \defaultfontfeatures[\rmfamily]{Ligatures=TeX,Scale=1}
\fi
\usepackage{lmodern}
\ifPDFTeX\else
  % xetex/luatex font selection
\fi
% Use upquote if available, for straight quotes in verbatim environments
\IfFileExists{upquote.sty}{\usepackage{upquote}}{}
\IfFileExists{microtype.sty}{% use microtype if available
  \usepackage[]{microtype}
  \UseMicrotypeSet[protrusion]{basicmath} % disable protrusion for tt fonts
}{}
\makeatletter
\@ifundefined{KOMAClassName}{% if non-KOMA class
  \IfFileExists{parskip.sty}{%
    \usepackage{parskip}
  }{% else
    \setlength{\parindent}{0pt}
    \setlength{\parskip}{6pt plus 2pt minus 1pt}}
}{% if KOMA class
  \KOMAoptions{parskip=half}}
\makeatother
\usepackage{longtable,booktabs,array}
\newcounter{none} % for unnumbered tables
\usepackage{calc} % for calculating minipage widths
% Correct order of tables after \paragraph or \subparagraph
\usepackage{etoolbox}
\makeatletter
\patchcmd\longtable{\par}{\if@noskipsec\mbox{}\fi\par}{}{}
\makeatother
% Allow footnotes in longtable head/foot
\IfFileExists{footnotehyper.sty}{\usepackage{footnotehyper}}{\usepackage{footnote}}
\makesavenoteenv{longtable}
\setlength{\emergencystretch}{3em} % prevent overfull lines
\providecommand{\tightlist}{%
  \setlength{\itemsep}{0pt}\setlength{\parskip}{0pt}}
\usepackage{bookmark}
\IfFileExists{xurl.sty}{\usepackage{xurl}}{} % add URL line breaks if available
\urlstyle{same}
% fallback for those not using the hyperref driver hyperxmp:
\makeatletter
\@ifundefined{xmpquote}{\newcommand{\xmpquote}[1]{#1}}{}
\makeatother
\hypersetup{
  pdftitle={Structural Preconditions of Anti-Systemic Political Emergence},
  pdfauthor={Satchel Hamilton · Civilodynamics Initiative},
  hidelinks,
  pdfcreator={LaTeX via pandoc}}

\title{Structural Preconditions of Anti-Systemic Political Emergence}
\usepackage{etoolbox}
\makeatletter
\providecommand{\subtitle}[1]{% add subtitle to \maketitle
  \apptocmd{\@title}{\par {\large #1 \par}}{}{}
}
\makeatother
\subtitle{A Civilodynamic Analysis of the United States, 2000--2024
(Working Paper No.~1, revised)}
\author{Satchel Hamilton · Civilodynamics Initiative}
\date{Canonical Corpus, Document 5 of 8 · Working Paper No.~1, Revision
2.0 · June 2026}

\begin{document}
\maketitle

\begin{quote}
\textbf{Status and role.} This is the flagship empirical application of
Civilodynamics: a structural analysis of the United States over
2000--2024 and the first registered set of prospective predictions. It
is a \emph{revision} of Working Paper No.~1 (April 2026). The
theoretical framework and all mathematics are specified in
\textbf{Documents 1, 2, and 4}, to which this paper defers; this paper
applies them. Three things in this revision differ materially from the
original and are stated up front because they bear on how the paper
should be read:

\begin{enumerate}
\def\labelenumi{\arabic{enumi}.}
\tightlist
\item
  \textbf{Every Civilizational Stress Index value here is an
  illustrative estimate, not a computed quantity.} The measurement
  pipeline (Document 4) has not yet been run end-to-end on these data.
  The numbers are indicative of magnitude and trajectory; they will
  change when computed. What is under test is the \emph{qualitative
  structural claim} and the \emph{registered predictions}, not the
  specific decimals.
\item
  \textbf{The interpretation of the stress index is corrected.} The
  original read a multivariate distance against the univariate 2σ/3σ
  rule, which is a mathematical error (Document 2 §4.4). The corrected
  reading is used throughout.
\item
  \textbf{The contested causal and political claims are demoted from
  assertions to explicitly operationalized, falsifiable sub-hypotheses},
  each with a stated way to be wrong (§4), and the paper's posture
  toward candidate-centered and contingency-based alternative
  explanations is even-handed.
\end{enumerate}

The five prospective predictions (§7) were registered in April 2026 and
are carried forward \textbf{unchanged}, retaining their original
timestamp, in keeping with the discipline's commitment not to revise
registered predictions retroactively.
\end{quote}

\begin{center}\rule{0.5\linewidth}{0.5pt}\end{center}

\subsection{Abstract}\label{abstract}

We apply Civilodynamics --- a framework that models societies as complex
dynamical systems and treats civilizational stress as the primary
structural variable behind political outcomes --- to the United States
over 2000--2024. The framework represents a society's stress state as a
vector across six blocks (economic, elite, legitimacy, cohesion,
epistemic, demographic) reduced to a scalar Civilizational Stress Index
(CSI). Using public-domain data (World Inequality Database, Bureau of
Labor Statistics, Pew, Gallup, Federal Reserve, and associated sources),
we present an \emph{illustrative} CSI trajectory and argue that the six
stress blocks were jointly elevated and rising \textbf{before} the
political developments of 2015--2016 that candidate-centered accounts
treat as causes. We advance, as the central hypothesis, that
anti-systemic political candidacy achieving major-party success is a
structurally probable output of high-stress states rather than a
candidate-specific anomaly --- and we state precisely how this
hypothesis, and three more contested sub-hypotheses it generates, could
be falsified. We develop a structural comparison with the Weimar
Republic (1919--1933), framed explicitly as
\textbf{hypothesis-generating analogy, not validation}, with its
disanalogies given equal weight. We register five timestamped,
Brier-scorable predictions for 2026--2030. All CSI values are
illustrative pending the computation specified in Document 4; the
substance under test is the qualitative structural account and the
registered predictions. We do not dispute that candidate-specific and
contingent factors influenced outcomes; we argue for structural
\emph{primacy} as a testable claim, and we treat the registered
predictions as its arbiter.

\textbf{Keywords:} civilodynamics, structural-demographic theory,
political instability, anti-systemic politics, elite overproduction,
epistemic coherence, Civilizational Stress Index, poliform transition,
U.S. political economy.

\begin{center}\rule{0.5\linewidth}{0.5pt}\end{center}

\subsection{1. Introduction}\label{introduction}

The U.S. electoral outcomes of 2016 and 2024 generated an extensive
literature of candidate-centered explanation: the liabilities of
specific nominees, late-breaking information events, the Electoral
College, polling error. These accounts are not wrong, and this paper
does not claim they are. It asks a different and complementary question:
were the conditions that produced these outcomes measurable in advance,
and do they constitute a structurally probable \emph{class} of
civilizational behavior rather than a sequence of contingent accidents?

This paper argues for the \textbf{structural-primacy hypothesis}: that
in high-stress civilizational states, the structural demand for
anti-systemic political expression rises substantially, that this demand
exists in the electorate prior to and independently of any specific
candidate, and that the candidate who channels it is a secondary
contingency relative to the structural condition. This is offered as a
hypothesis with stated falsification conditions (§4), not as an
established result, and it is held alongside --- not against ---
candidate-centered explanation: structural conditions set the
possibility space; candidate factors operate within it and shape
margins.

The most proximate antecedent is the structural-demographic theory of
Turchin and colleagues (Turchin 2003, 2016; Turchin and Nefedov 2009),
which identifies elite overproduction, popular immiseration, and state
fiscal stress as primary drivers of political instability.
Civilodynamics extends this in the three narrow ways set out in the
charter (Document 1): an explicit epistemic/information layer, the
poliform--morphospace integration, and a real-time measurement ambition.
The framework's debt to cliodynamics is foundational and is acknowledged
as such.

In the poliform language of Document 3, the United States is an instance
of the \textbf{industrial nation-state} basin, and the sharp structural
question this paper poses (§6) is whether the elevated stress of
2008--2024 is an \emph{intra-basin} secular-cycle peak that will relax,
or the early signature of a \emph{between-basin} transition toward the
algorithmic-surveillance basin. The paper does not resolve this
question; it frames it and registers predictions that bear on it.

The paper proceeds: §2 states the framework briefly; §3 presents the
component-level stress trajectory and an illustrative CSI; §4 gives the
retrospective interpretation as a set of falsifiable sub-hypotheses,
with alternatives; §5 develops the Weimar comparison as
hypothesis-generating analogy; §6 poses the poliform question; §7
registers five predictions; §8 states limitations; §9 concludes.

\subsection{2. Framework (in brief)}\label{framework-in-brief}

The full specification is in Documents 1, 2, and 4. The minimum needed
here:

\textbf{The stress readout and the CSI.} A society's stress state is a
vector \(\mathbf{u}(t)\) across six blocks --- economic stress, elite
dynamics, legitimacy, cohesion, epistemic coherence, demographic
pressure --- each built from operationalized variables drawn from public
data (the full variable table and sources are in Document 4 §3). The CSI
is a scalar reduction of \(\mathbf{u}\) measuring distance from a
defended stable baseline.

\textbf{The baseline and its contestability.} Variables are normalized
against a U.S. stable-period baseline of 1948--1973 (post-war
convergence: falling inequality, broadly rising real wages,
comparatively high institutional trust). This baseline is contestable,
and the paper acknowledges it directly: the same period was marked by
racial apartheid, Cold War anxiety, and gendered institutional
exclusion. The choice is a methodological commitment subject to
sensitivity analysis (§8), not a free parameter, and the
baseline-covariance problem it creates is resolved per Document 4 §4 (a
coverage-restricted baseline vector plus a modern-era epistemic
extension).

\textbf{The CSI's correct interpretation.} The CSI is reported here as a
per-dimension-normalized composite (Document 2 §4.5), not as a raw
Mahalanobis distance read against the 2σ/3σ rule. This matters: a raw
multivariate Mahalanobis distance is chi-squared distributed and a
baseline-typical state already sits well above zero, so the original
threshold labels were miscalibrated. The corrected reading is
qualitative-magnitude-preserving --- the trajectory still rises --- but
the specific numbers and their thresholds await the computation of
Document 4.

\textbf{Anti-Systemic Political Event (ASPE).} Defined in advance to
avoid post-hoc classification, and deliberately broad: (a) a candidate
without prior major-party institutional career achieving major-party
presidential nomination; (b) a candidate campaigning on institutional
delegitimation exceeding 35\% of the two-party popular vote; or (c) a
third-party/independent candidacy exceeding 15\% of the popular vote.
The core hypothesis treats elevated CSI as a \textbf{necessary but not
sufficient} condition for ASPE occurrence --- a claim that, as §4 notes,
requires a stated functional form and base rate to be genuinely
falsifiable, which the predictions of §7 supply.

\subsection{3. The U.S. stress trajectory,
2000--2024}\label{the-u.s.-stress-trajectory-20002024}

The component-level figures below are largely \textbf{sourced reported
statistics} (and should be verified against current provider data at
submission); their normalization and aggregation into a CSI is
\textbf{illustrative/estimated} pending Document 4's pipeline.

\textbf{Economic stress.} Real wages for the bottom two quintiles were
effectively flat from 1979 through approximately 2015 even as
productivity rose (Mishel et al.~2015; EPI wage-productivity series).
Top-1\% net wealth share rose from roughly 23\% (1980) toward roughly
38\% (2020) per the Federal Reserve DFA, approaching late-1920s levels
(Saez and Zucman 2020), corroborated by WID. Intergenerational mobility
declined over the period (Chetty et al.~2014). \emph{(Specific
transition rates cited in the original --- e.g., bottom-to-top-quintile
probabilities --- should be re-verified against the Opportunity Insights
series before submission; they are not load-bearing for the
trajectory.)}

\textbf{Elite dynamics.} The Elite Overproduction Index (construction
and caveats in Document 4 §6) rose materially over the period.
Law-school first-year enrollment rose to a peak near 52,500 around 2010
(ABA) while lawyer-position growth lagged; PhD conferrals rose toward
roughly 55,000/year by 2010 (NCES/SED) while tenure-track availability
in the humanities and social sciences contracted after
\textasciitilde2008, widening the credential-position gap. Turchin
(2016) identifies elite overproduction as among the most predictive
instability variables in his historical set. \emph{(EOI values are
estimated pending the harmonized protocol of Document 4 §6, which must
bridge the 2000/2010 occupational reclassifications.)}

\textbf{Legitimacy.} Gallup's institutional-confidence series shows mean
confidence across consistently-tracked institutions falling from roughly
48\% (mid-1970s) toward roughly 27\% (2023); Congressional confidence
specifically fell to single digits (\textasciitilde8\%, 2023). Pew's
parallel series shows accelerating partisan divergence in trust after
\textasciitilde2008.

\textbf{Cohesion.} Population-weighted county partisan lean increased
over the period (Bishop 2008; Chen and Rodden 2013), reducing routine
cross-partisan contact; affective polarization rose and, per Iyengar et
al.~(2019), exceeds issue-based polarization and accelerated after
\textasciitilde2000.

\textbf{Epistemic coherence.} This is the weakest-measured block
(Document 4 §5). The Media Overlap Score, estimable with quality only
from \textasciitilde2012 and reconstructed earlier, declined over the
period; the Factual Consensus Rate \textbf{has no instrument yet} and
its pre-instrument values are not reported as data. The composite ECI is
estimated to have moved from a functional into a degraded range over
\textasciitilde2012--2014, but this estimate rests on interpolation and
is explicitly indicative, not measured.

\textbf{Demographic.} Youth-cohort and rural-urban divergence measures
(Census ACS; ERS) contributed secondary, slower-moving pressure.

\subsubsection{The illustrative CSI
trajectory}\label{the-illustrative-csi-trajectory}

The table below is \textbf{illustrative}. It expresses an estimated
per-dimension-normalized composite; the regime bands are illustrative
ranges on that normalized scale, \textbf{not} thresholds derived from
the corrected χ²/Hotelling calibration (which awaits computation,
Document 4 §4). It is included to show the shape of the claim, not to
assert measured values.

{\def\LTcaptype{none} % do not increment counter
\begin{longtable}[]{@{}
  >{\raggedright\arraybackslash}p{(\linewidth - 6\tabcolsep) * \real{0.0759}}
  >{\raggedright\arraybackslash}p{(\linewidth - 6\tabcolsep) * \real{0.3797}}
  >{\raggedright\arraybackslash}p{(\linewidth - 6\tabcolsep) * \real{0.2658}}
  >{\raggedright\arraybackslash}p{(\linewidth - 6\tabcolsep) * \real{0.2785}}@{}}
\toprule\noalign{}
\begin{minipage}[b]{\linewidth}\raggedright
Year
\end{minipage} & \begin{minipage}[b]{\linewidth}\raggedright
Illustrative normalized CSI
\end{minipage} & \begin{minipage}[b]{\linewidth}\raggedright
Illustrative regime
\end{minipage} & \begin{minipage}[b]{\linewidth}\raggedright
Primary contributors
\end{minipage} \\
\midrule\noalign{}
\endhead
\bottomrule\noalign{}
\endlastfoot
2000 & \textasciitilde1.2 & Moderate & Wealth concentration, wage
stagnation \\
2004 & \textasciitilde1.4 & Moderate & Legitimacy decline, EOI rising \\
2008 & \textasciitilde1.8 & Elevated & Financial-crisis amplification \\
2010 & \textasciitilde2.0 & Elevated & EOI peak, legitimacy decline \\
2012 & \textasciitilde2.2 & High & Epistemic deterioration onset,
sorting \\
2015 & \textasciitilde2.5 & High & All blocks elevated \\
2016 & \textasciitilde2.6 & High & Election year \\
2020 & \textasciitilde2.9 & Critical & Pandemic amplification \\
2022 & \textasciitilde2.8 & Critical & Legitimacy decline deepened \\
2024 & \textasciitilde3.0 & Critical & All blocks at/near period
maxima \\
\end{longtable}
}

\textbf{The load-bearing claim is not any number in this table.} It is
the \emph{precedence} claim: the six stress blocks were jointly elevated
and rising \textbf{before} the 2015--2016 political developments. That
claim rests on the component evidence (wage stagnation, wealth
concentration, rising elite overproduction, legitimacy collapse, early
epistemic fragmentation), most of which is reasonably solid and
pre-2015, and it is falsifiable: if computation shows the blocks were
\emph{not} jointly elevated before 2015, the structural-primacy
hypothesis loses its principal support.

\subsection{4. Retrospective interpretation: 2016 and 2024 as a set of
testable
hypotheses}\label{retrospective-interpretation-2016-and-2024-as-a-set-of-testable-hypotheses}

This section states the structural interpretation as falsifiable
sub-hypotheses rather than assertions, gives the alternatives equal
hearing, and ties each hypothesis to the registered predictions that
test it.

\subsubsection{H1 --- Structural precedence (the core
claim)}\label{h1-structural-precedence-the-core-claim}

\emph{Statement:} Stress was elevated and rising before the 2015--2016
developments, such that anti-systemic emergence is better explained as a
structurally probable output of a high-stress state than as a
candidate-specific anomaly. \emph{Operationalization:} the six stress
blocks jointly above baseline and trending upward by 2012--2015.
\emph{Way to be wrong:} computation (Document 4) showing the blocks were
not jointly elevated/rising before 2015 falsifies H1. \emph{Supporting
context, stated even-handedly:} the anti-establishment signal in 2016
was visible across multiple candidates of both parties before any one
became dominant, and broadly contemporaneous anti-establishment outcomes
appeared in several democracies (Norris and Inglehart 2019) ---
consistent with a common structural driver. \emph{But the cross-national
pattern is not uniform:} several democracies under comparable structural
stress did \textbf{not} produce equivalent anti-systemic outcomes in the
same window, which is itself evidence that structure sets a possibility
space rather than determining outcomes. H1 claims primacy of structure,
not sufficiency.

\subsubsection{H2 --- Centrist structural
disadvantage}\label{h2-centrist-structural-disadvantage}

\emph{Statement:} In high-CSI environments, institutionalist
(incremental-reform, pro-institution) candidacies face a structural
headwind independent of the personal qualities of the specific
candidate. \emph{Operationalization:} the vote share of institutionalist
candidacies declines as CSI rises, controlling for candidate-quality
proxies, across cases. \emph{Way to be wrong (stated in the original and
retained):} an institutionalist candidate achieving a decisive
popular-vote victory in a CSI-high environment falsifies H2. \emph{This
is directly tested by Prediction 1.} The institutionalist nominees of
2016 and 2024 (Clinton; Harris) are the U.S. instances; H2 makes a claim
about the \emph{environment}, not about those individuals' competence or
policy positions, which it treats as secondary.

\subsubsection{H3 --- Delegitimation-signal
inversion}\label{h3-delegitimation-signal-inversion}

\emph{Statement:} In collapsed-legitimacy environments, signals of
institutional disapproval (indictments, elite condemnation,
mainstream-media opposition) can function as positive credentials of
anti-systemic authenticity for a segment of the electorate.
\emph{Operationalization:} among voters with low institutional trust,
exposure to institutional-disapproval signals is associated with
non-negative or positive candidate evaluation, versus negative
association among high-trust voters. \emph{Way to be wrong:} if, across
measurement, institutional-disapproval signals depress support uniformly
regardless of institutional-trust level, H3 fails. \emph{Caveats given
equal weight:} the counterfactual is unobservable (we cannot run the
2024 election without the indictments), the net effect may be mixed and
masked by other factors, and the claim is among the most contested in
the framework. The paper states it as a hypothesis precisely because it
is not established. \textbf{This is a structural description of a
hypothesized mechanism and carries no normative endorsement of any
candidate, outcome, or conduct.}

\subsubsection{The information layer --- mechanism, not
intent}\label{the-information-layer-mechanism-not-intent}

The framework holds that engagement-optimizing information environments
amplify emotionally arousing, identity-confirming, low-accuracy content,
degrading the shared factual commons and increasing the electoral
fitness of candidates optimized for algorithmically-amplified engagement
over institutional competence. This \textbf{mechanism} is
operationalizable (via the ECI, Document 4 §5) and is retained. The
paper deliberately does \textbf{not} advance the stronger claim that
this degradation was \emph{coordinated or deliberately engineered}: that
is a causal-intent claim that is far harder to establish, pulls toward
unfalsifiability, and is not needed for the structural argument. The
mechanism is the claim; intent is not asserted.

\subsubsection{Alternatives the paper does not
dispute}\label{alternatives-the-paper-does-not-dispute}

Candidate-specific factors (nominee liabilities, communication styles),
contingent events (information shocks, the pandemic's timing), and
institutional features (the Electoral College, primary dynamics) all
influenced outcomes and margins. The structural account claims
\emph{primacy} --- that these operated within a possibility space the
structural conditions defined --- not \emph{exclusivity}. Distinguishing
structural primacy from these alternatives is exactly what the
registered predictions (§7) are designed to do over the coming cycles.

\subsection{5. Historical comparison: the Weimar Republic,
1919--1933}\label{historical-comparison-the-weimar-republic-19191933}

\begin{quote}
\textbf{Epistemic status --- read this first.} The Weimar comparison is
\textbf{hypothesis-generating analogy, not validation.} It is a sample
of one, scored qualitatively. It cannot validate the CSI scale, because
no CSI has been computed for Weimar on commensurable variables; doing so
is deferred to the Retrodiction Casebook (Document 6). The original
framing of this comparison as ``a form of cross-civilizational
validation'' is withdrawn. The parallels suggest hypotheses; the
disanalogies are given equal weight; nothing here implies the United
States is on the Weimar trajectory.
\end{quote}

\textbf{Structural parallels (qualitative).} Across economic stress
(sustained immiseration and wealth divergence), elite overproduction (a
frustrated credentialed cohort), institutional legitimacy (a
delegitimated parliament/Congress), and epistemic fragmentation (a
party-press commons in Weimar; an algorithmic one now), the two cases
show directionally similar structural conditions. The \emph{mechanisms}
differ sharply --- a partisan press versus algorithmic curation --- even
where the structural outcome (a fragmented factual commons) rhymes.

\textbf{Critical disanalogies (given equal weight).} The Weimar
constitution had specific structural weaknesses (Article 48 emergency
powers; no constructive vote of no confidence) that the U.S. system does
not share in the same form; U.S. federalism, an independent judiciary,
and civil-society density provide resistance Weimar lacked. The German
military's institutional ambivalence toward democracy has no direct
contemporary U.S. parallel. The Weimar monetary catastrophe
(hyperinflation then deflation) has no U.S. counterpart. And Weimar's
organized paramilitary mobilization (SA, SS, Reichsbanner) at scale is
not present in the contemporary U.S. case, though January 6, 2021 is a
qualitatively novel U.S. data point warranting structural monitoring.
These disanalogies justify the conclusion that the U.S. system, while in
a high-stress state, is \textbf{not} in Weimar's terminal trajectory ---
and they justify treating the comparison as a source of hypotheses, not
a forecast.

\subsection{6. The poliform question}\label{the-poliform-question}

Translating into the poliform framework (Document 3): the United States
is in the \textbf{industrial nation-state} basin. The central structural
question is whether the elevated stress of 2008--2024 is (a) an
\textbf{intra-basin secular-cycle peak} --- a Turchin ``Age of Discord''
that, however severe, relaxes within the basin --- or (b) the early
phase of a \textbf{between-basin transition} toward the
algorithmic-surveillance basin, driven by a regime change in the
information axis (a Type-2 basin-reshaping in the terms of Document 2
§6.3).

Two honest caveats sharpen this. First, the industrial nation-state is
the \emph{weakest convergence case} in the catalog (Document 3 §5): it
arose from few origins and spread by diffusion, so claims about its
``natural'' instabilities are less secure than for, say, the
steppe-nomad confederacy. Second, the algorithmic-surveillance basin is
\emph{untested} --- possibly a genuine new attractor, possibly a
transient, possibly an artifact of shared technology diffusing across
different political traditions. The framework does not claim to know
which. It claims only that the question is now answerable in principle,
and the registered predictions bear on it: sustained ECI decline and
continued anti-systemic consolidation would be consistent with (b);
relaxation of the stress blocks would be consistent with (a).

\subsection{7. Prospective predictions,
2026--2030}\label{prospective-predictions-20262030}

\begin{quote}
Registered April 2026; carried forward \textbf{unchanged} in this
revision, retaining the original timestamp. Resolution criteria are
clarified for precision where needed but the claims and probabilities
are not altered. The probabilities are \textbf{calibrated subjective
estimates} informed by the structural assessment and by analogy to
historical cases --- they are not outputs of a fitted model (no model is
yet fitted), and they therefore carry model uncertainty. All are
Brier-scored at resolution; updated estimates, if any, will be published
as separately-timestamped versions.
\end{quote}

\textbf{Prediction 1 --- P = 0.72.} At least one major-party
presidential or vice-presidential nomination in the 2028 cycle goes to a
candidate whose platform includes explicit institutional delegitimation
of at least one co-equal branch of government. \emph{Resolution (Nov
2028):} documented campaign statements of a nominated candidate
explicitly calling for subordination, defunding, restructuring, or
delegitimation of the judiciary, legislature, or executive \emph{as an
institution} (not mere policy disagreement with current officeholders).
\emph{Tests H1 and H2.}

\textbf{Prediction 2 --- P = 0.78.} The U.S. ECI, measured by the
validated methodology of Document 4, shows continued decline from its
2024 level through 2028 absent a major intervention in platform
incentive architecture. \emph{Resolution (Dec 2028):} ECI(2028)
\textless{} ECI(2024) under the standardized methodology.
\emph{Conditional on the ECI instrument being fielded; tests the
information-layer mechanism.}

\textbf{Prediction 3 --- P = 0.65.} At least three U.S. states
experience significant fiscal-stress events (general-obligation credit
downgrade, declared fiscal emergency, or failure to meet scheduled
pension obligations) between 2026 and 2030. \emph{Resolution (Dec
2030):} documented occurrence in ≥3 distinct states.

\textbf{Prediction 4 --- P = 0.42.} A significant anti-systemic
political formation distinct from the current Republican Party achieves
double-digit popular-vote share in at least one statewide or federal
election by 2030. \emph{Resolution (Nov 2030):} any non-Democratic,
non-Republican candidate exceeding 10\% in a Senate, gubernatorial, or
presidential election. \emph{The lower probability reflects genuine
uncertainty about whether the two-party structure absorbs anti-systemic
demand.}

\textbf{Prediction 5 --- P = 0.83.} The U.S. CSI does not return to a
moderate-stress regime by 2030 absent a structural policy intervention
sufficient to reverse at least two of: wealth-concentration trajectory,
bottom-quintile real-wage growth, institutional trust, or the EOI.
\emph{Resolution (Dec 2030):} assessed against the computed CSI under
Document 4's methodology. \emph{The most structurally grounded
prediction, resting on the slowness of these variables relative to
political cycles.}

A note on reflexivity (Document 2 §8.1): publishing these predictions
may influence the actors they concern. The discipline treats this as
unavoidable and addresses it through transparency; notably, an outcome
that falsifies Prediction 5 --- a structural correction --- is one the
author would welcome.

\subsection{8. Limitations and future
research}\label{limitations-and-future-research}

\textbf{Estimated, not computed.} The principal limitation is the one
stated throughout: the CSI here is illustrative. Converting it to a
computed result requires Document 4's pipeline --- the
coverage-restricted baseline, the corrected χ²/Hotelling thresholds, the
per-dimension-normalized headline, and uncertainty propagation. Until
then, this paper is a structured argument with indicative magnitudes,
not a measurement.

\textbf{Epistemic-layer measurement.} The ECI is the weakest block: one
component is reconstructed/estimated before \textasciitilde2012, the
other has no instrument yet. Prediction 2 is conditional on building it
--- and the instrument it requires (the FCR) is now specified and
pre-registered in \textbf{Document 9}, so fielding it is what makes
Prediction 2 live. This is the program's highest-priority empirical
investment (Document 4 §5).

\textbf{Baseline sensitivity.} All z-scores depend on the 1948--1973
baseline. Sensitivity analyses against 1948--1968 and 1950--1980
baselines are required; the original reported these as broadly
shape-preserving but level-shifting by \textasciitilde15--25\%,
indicating the trajectory is robust to the choice while specific
thresholds are not.

\textbf{The convergence-versus-diffusion caveat.} Because the industrial
nation-state is a weak convergence case (Document 3 §5), structural
claims that lean on its being a well-attested attractor with
characteristic instabilities should be appropriately hedged.

\textbf{Sample-of-one comparison.} Weimar generates hypotheses; it does
not validate. Genuine cross-civilizational validation requires the
computed retrodiction set of Document 6 (late Roman Republic, French
Revolution, Soviet collapse, 1848), against the
80\%-correct-characterization bar of the charter.

\textbf{Reflexivity.} Treated as an unavoidable feature addressed by
transparency (§7).

\subsection{9. Conclusion}\label{conclusion}

This paper has applied Civilodynamics to the United States, 2000--2024,
and registered five falsifiable predictions. Its central claim is the
structural-primacy hypothesis: that the six stress blocks were jointly
elevated and rising before the 2015--2016 developments, that
anti-systemic emergence is therefore better read as a structurally
probable output of a high-stress state than as a candidate-specific
anomaly, and that this --- together with three more contested
sub-hypotheses --- can be stated precisely enough to be wrong.

The paper is careful about what it has and has not shown. The CSI values
are illustrative pending computation; the Weimar comparison generates
hypotheses rather than validating; the contested causal claims are
demoted to operationalized sub-hypotheses with stated falsifiers; and
candidate-centered and contingent explanations are credited as operating
within, not against, the structural account. What distinguishes the
approach is not the confidence of its numbers --- which are provisional
--- but the structure of its commitments: a measurable stress state, a
falsifiable set of hypotheses, and a registered, Brier-scorable
predictive record that the coming cycles will test. Whether the specific
predictions prove right or wrong, that structure is what separates this
from event-level commentary, and the resolution dates of §7 are when the
framework's first real empirical verdict arrives.

\begin{center}\rule{0.5\linewidth}{0.5pt}\end{center}

\subsection{References}\label{references}

American Bar Association (2023); Bishop, B. (2008); Bureau of Labor
Statistics (2024, CES and OEWS); Chen, J., \& Rodden, J. (2013); Chetty,
R., Hendren, N., Kline, P., \& Saez, E. (2014); Congressional Budget
Office (2024); Dell, F., Piketty, T., \& Saez, E. (2007); Evans, R. J.
(2003); Federal Reserve Board (2024, DFA); Fritzsche, P. (1998); Gallup
Organization (2023); International Monetary Fund (2023, HPDD); Iyengar,
S., Lelkes, Y., Levendusky, M., Malhotra, N., \& Westwood, S. J. (2019);
James, H. (1986); Kahan, D. M. (2017); Kershaw, I. (1998); MIT Election
Data and Science Lab (2023); Mishel, L., Gould, E., \& Bivens, J.
(2015); National Center for Education Statistics (2022); Norris, P., \&
Inglehart, R. (2019); Pennycook, G., \& Rand, D. G. (2019); Peukert, D.
J. K. (1989); Pew Research Center (2022, 2023); Piketty, T. (2014);
Putnam, R. D. (2000); Reuters Institute (2023); Saez, E., \& Zucman, G.
(2020); Turchin, P. (2003, 2016); Turchin, P., \& Nefedov, S. A. (2009);
World Inequality Database (2024).

\emph{Bullock, J. G., et al.~(2015) and Prior, M., et al.~(2015), on
partisan expressive responding and incentivized-accuracy designs, are
cited in support of the ECI/FCR methodology in \textbf{Document 9} (the
FCR instrument specification and pre-registration).}

\emph{Working paper, revised. Cite with caution. All CSI values are
illustrative pending computation; all component statistics should be
verified against current provider data before submission. Framework:
Documents 1, 2, 4. Validation: Document 6.}

\end{document}
